% ----------------------------------------------------
% Projeto: UnDF-TeX
% Arquivo: 04-analise-e-discussao.tex
%
% Copyright (c) 2026
% Universidade do Distrito Federal (UnDF)
%
% Este arquivo faz parte do projeto UnDF-TeX.
%
% Licenciado sob a LaTeX Project Public License (LPPL),
% versão 1.3c ou posterior.
%
% Implementação técnica baseada nas normas:
% ABNT NBR 14724:2024
% ABNT NBR 10520:2023
%
% Este projeto não reproduz o texto integral
% das normas da ABNT.
% ----------------------------------------------------

\chapter{Análise e discussão}

Este capítulo reúne exemplos práticos de elementos frequentemente presentes em um TCC.

\section{Exemplo de equação}

A equação de segundo grau pode ser representada por:
\begin{equation}
	x = \frac{-b \pm \sqrt{b^2 - 4ac}}{2a}
\end{equation}

Também é possível usar uma fórmula destacada sem numeração:
\[
E = mc^2
\]

\section{Exemplo de tabela}

\begin{table}[htbp]
	\centering
	\caption{Etapas do planejamento do TCC}
	\begin{tabular}{p{0.28\textwidth} p{0.45\textwidth}}
		\hline
		\textbf{Etapa} & \textbf{Finalidade} \\
		\hline
		Tema & Definir o foco do trabalho \\
		Problema & Formular a questão central \\
		Objetivo & Delimitar o que se pretende alcançar \\
		\hline
	\end{tabular}
\end{table}

A Tabela 4.1 sintetiza as etapas iniciais do planejamento do TCC.

\section{Exemplo de figura}
\undffigura[0.50\textwidth]{figuras/exemplo.jpg}{The Coppersmith}{ilustracao-tcc}[Edgar Melville Ward]

Como mostra a \figuraref{ilustracao-tcc}, a imagem pode ser utilizada como recurso visual de apoio ao texto acadêmico.

\section{Exemplo de figura com fonte externa}
\undffigura[0.50\textwidth]{figuras/exemplo2.jpg}{The Starry Night}{ilustracao-fonte-externa}[Vincent Van Gogh, 1889]

\section{Exemplo de articulação interna}

O tema foi apresentado anteriormente na Seção 1.1. O instrumento complementar está no Apêndice A. O documento de apoio está no Anexo A.
